% ==========================================================================
% GENERATED by scripts/split_paper.py. Do not edit this file.
%
% Source: imported_papers/grokking/iclr2027_conference_v2.tex
% Part 10: Interventions that separate credit from the gradient
%
% Edits belong upstream, in the published paper. To change what the
% chapter includes, change imports/grokking.toml and re-run
% the splitter.
% ==========================================================================
\section{Interventions that separate credit from the gradient}
Section~\ref{sec:law} shows that the circuit a rule builds matches the credit it applies. That match invites one objection: a unit's incoming gradient norm is proportional to its outgoing weight norm times its backpropagated error, so credit and circuit are coupled by arithmetic, and their agreement could be forced rather than causal. This section breaks the coupling from three sides, varying the rule, then the credit, then the weights.

\subsection{Published per-unit rules and a weight geometry control}
\label{app:named}

Nero \citep{liu2021nero} normalizes the update per neuron, which is the per-unit grain, so the theory says it must land distributed. It does, at two learning rates (Table~\ref{tab:named}), and its tolerance to targeted deletion, $0.626$, is the highest of any rule we measure.

If the knob worked only by keeping the input rows at equal norm, the result would be a fact about weight shape rather than credit. Forcing each row to a fixed norm after every step, with no rescaling of the gradient, leaves the circuit near backpropagation, at participation $0.661\,(0.015)$ against $0.595\,(0.023)$, with lower deletion tolerance (Table~\ref{tab:named}). Equal row norms alone do not do what the credit rescaling does. Weight normalization \citep{salimans2016weightnorm} extends that control to a published method. It reparameterizes each unit's incoming vector as one learned scale times a direction, a per-unit constraint on the weights that rescales no credit, and under plain backpropagation, with weight decay on the per-unit scales, it groks to test $0.993$ and stays concentrated at participation $0.466\,(0.026)$ over three seeds. A ubiquitous per-unit reparameterization of the weights lands with the row-norm control and not with Nero, which is the dissociation between weight geometry and credit at publication scale.

The sign-SGD sweep guards a boundary rather than the coupling: Table~\ref{tab:twoaxis} uses sign-SGD as the case of uniform credit that still memorizes, which only holds if the failure is not a tuning artifact. At learning rates $10^{-4}$ and $10^{-5}$ and at weight decay $10^{-2}$, it reaches test $0.000$ in every case while delivering credit participation $1.000$, so it is not. The failure also survives gating: multiplying the sign update by the current loss, so the step vanishes as the fit is reached, still memorizes at every learning rate from $10^{-4}$ to $10^{-2}$ over three seeds, which keeps the sign family scored by its convergence theory rather than by the fixed point case of Proposition~\ref{prop:persist}.

\subsection{Imposed credit profiles}
\label{app:imposed}
Imposed profiles are the direct test of the coupling objection, since the credit is supplied from outside the gradient and nothing the network computes can feed back into it. At every step we replace the per-unit magnitudes of the input layer update with a profile drawn to a chosen participation, keeping each row's direction and the summed norm, and dealing the profile to units at random so it carries no information about any unit's role. The grokked circuit follows the imposed value across the sweep, and deletion tolerance follows with it, at test accuracy above $0.96$ everywhere (Table~\ref{tab:named}). A profile uncorrelated with the gradient leaves no coupling to explain the match. Participation is a quantity we can set, and the circuit obeys it.

The last row is the sharpest test. It keeps backpropagation's own credit magnitudes and only changes who receives them, dealing them out in a fresh random order every step. The set of magnitudes is unchanged, so participation is identical at every step, and only persistence is destroyed: the unit receiving the largest share now is not the one receiving it next. That alone moves the circuit to $0.992$ against backpropagation's $0.613$. A concentrated circuit needs credit that concentrates on the \emph{same} units over time, since concentration that moves around averages away.

Two boundary notes. Concentrating credit does not undo a circuit that is already distributed: rescaling DFC's magnitudes away from uniform drives credit participation from $0.92$ to $0.31$ and per-neuron periodicity from $0.89$ to $0.33$, while test accuracy and the frequency spread stay put, so per-neuron structure follows the credit applied now and frequency coverage, once built, persists. And nothing here depends on initialization scale: across scales of $0.5$ to $4$, backpropagation stays between $0.586$ and $0.613$, the knob between $0.990$ and $0.998$, and applied credit tracks the grokked circuit to a mean absolute deviation of $0.008$ throughout.

The knob is also local to the layer it acts on. Applied to the second hidden layer instead of the first, that layer's circuit distributes to $1.000$ while the first layer's stays at plain flow, $0.593\,(0.004)$ with $86$ units idle, and applied to both layers it distributes both, at test accuracy $0.98$ and above throughout (three seeds). Each layer's circuit follows its own units' credit, so the agreement between credit and circuit is not one layer's norms dragging another's along.

\begin{table}[!ht]
\caption{Nero lands at the distributed end as its per-neuron grain predicts, equal row norms without credit rescaling stay near backpropagation, and the grokked circuit follows an imposed credit profile to a mean absolute deviation of $0.018$ even though the profile carries no information about any unit's role. Interventions that vary the rule and interventions that vary the credit ($p=31$, three seeds, standard deviations in parentheses, "no" meaning the rule does not grok). In the upper block the row-norm constraint keeps row norms equal without rescaling credit, and sign-SGD fails at every setting swept. In the lower block each profile is drawn to the target participation and assigned to units at random, and the last row keeps backpropagation's own magnitudes and permutes which unit receives which at every step. Each block has its own backpropagation run.}
\label{tab:named}
\begin{center}
\small
\setlength{\tabcolsep}{4pt}
\begin{tabular}{lccccc}
\toprule
Rule or condition & Test & Credit PR & Circuit PR & Random del & Targeted del \\
\midrule
\multicolumn{6}{l}{\emph{Varying the rule}} \\
Backpropagation & $0.998$ & $0.586$ & $0.595\,(0.023)$ & $0.672$ & $0.311$ \\
Nero \citep{liu2021nero} & $0.985$ & $0.968$ & $0.986\,(0.002)$ & $0.748$ & $0.626$ \\
Nero, lr $3\times10^{-3}$ & $0.978$ & $0.959$ & $0.983\,(0.001)$ & $0.728$ & $0.541$ \\
Row-norm constraint & $0.954$ & $0.611$ & $0.661\,(0.015)$ & $0.607$ & $0.256$ \\
Sign-SGD, lr $10^{-4}$ & $0.000$ & $1.000$ & no & --- & --- \\
Sign-SGD, lr $10^{-5}$ & $0.000$ & $1.000$ & no & --- & --- \\
Sign-SGD, lr $10^{-5}$, wd $10^{-2}$ & $0.000$ & $1.000$ & no & --- & --- \\
\midrule
\multicolumn{6}{l}{\emph{Varying the credit}} \\
Backpropagation & $0.997$ & $0.604$ & $0.613\,(0.016)$ & $0.673$ & $0.298$ \\
Imposed $0.60$ & $0.980$ & $0.600$ & $0.637\,(0.003)$ & $0.607$ & $0.240$ \\
Imposed $0.70$ & $0.981$ & $0.700$ & $0.727\,(0.001)$ & $0.663$ & $0.319$ \\
Imposed $0.80$ & $0.972$ & $0.800$ & $0.815\,(0.005)$ & $0.693$ & $0.420$ \\
Imposed $0.90$ & $0.969$ & $0.900$ & $0.902\,(0.005)$ & $0.723$ & $0.533$ \\
Imposed $1.00$ & $0.975$ & $1.000$ & $0.990\,(0.003)$ & $0.775$ & $0.599$ \\
\midrule
Reshuffled each step & $0.974$ & $0.979$ & $0.992\,(0.002)$ & $0.775$ & $0.601$ \\
\bottomrule
\end{tabular}
\end{center}
\end{table}

\subsection{Participation as a sufficient statistic}
\label{app:suff}
The knob reaches its participation by rescaling credit. A second route reaches the same participation without touching the gradient: constrain the outgoing norms of the first hidden layer to a chosen profile after every step. Comparing the two at matched participation asks which consequences follow from participation itself and which need the credit rescaling (Table~\ref{tab:suff}).

\paragraph{Deletion tolerance follows participation alone.} Proposition~\ref{prop:fluct} says it should, and at participation $0.99$ the two routes both reach random deletion tolerance $0.777$, matching to three decimals. That measure is computed from the norms we matched on, so we also ran the one attack that never consults a norm, the greedy search for the smallest set of units whose deletion drops test below $0.5$. It needs $66.7\,(11.1)$ units against the knob and $68.7\,(10.8)$ against the constraint, the same within seed spread. Targeted deletion, by contrast, is not comparable here: the constrained profile at $0.99$ is exactly uniform, so the attack has no highest norm unit to remove and degenerates to random deletion, which is why the $0.747$ in the table cannot be compared with the knob's $0.590$.

\paragraph{Adaptation time mostly follows it too.} At matched participation the constrained network takes $1480$ epochs against the knob's $1608$, and within the constrained family adaptation rises with participation, $740$, $1228$ and $1480$ epochs at $0.60$, $0.80$ and $0.99$. The remainder is spare capacity: at participation $0.60$ backpropagation keeps $87$ free units and adapts in $464$ epochs where the constrained network, with none, takes $740$, the split Appendix~\ref{app:adapt} details.

\paragraph{Per-neuron periodicity does not follow it.} The constrained family is flat, $0.865$ to $0.880$ from participation $0.60$ to $0.99$, while backpropagation at the same $0.60$ sits at $0.667$. Two networks with identical participation differ in per-neuron structure by about $0.20$, so periodicity is not a function of participation, and its agreement with the knob at $0.99$ is only both ends being high. This bounds the claim: participation is close to sufficient for the damage and adaptation consequences, and not for how cleanly single units tune.

\begin{table}[!ht]
\caption{Matching participation by constraining outgoing norms reproduces the deletion tolerance, so participation carries that consequence. The constraint route to the same participation (modular addition, $p=31$, five seeds, standard deviations in parentheses). Adaptation is epochs to test $0.9$ on subtraction, measured in this five seed run, so the backpropagation and knob values differ from the ten seed numbers of Table~\ref{tab:adapt}. The targeted deletion column for the constrained rows is an artifact of the exactly uniform profile, as explained above.}
\label{tab:suff}
\begin{center}
\footnotesize
\setlength{\tabcolsep}{4pt}
\begin{tabular}{lcccccc}
\toprule
Condition & Test & Circuit PR & Random del & Targeted del & Periodicity & Adapt \\
\midrule
Backprop             & $0.995\,(0.007)$ & $0.597\,(0.020)$ & $0.669$ & $0.321$ & $0.667$ & $464$ \\
Knob ($\beta{=}1$)   & $0.976\,(0.015)$ & $0.993\,(0.002)$ & $0.777$ & $0.590$ & $0.884$ & $1608$ \\
\midrule
Constrained $0.60$   & $0.987\,(0.016)$ & $0.600\,(0.000)$ & $0.647$ & $0.357$ & $0.865$ & $740$ \\
Constrained $0.80$   & $0.977\,(0.008)$ & $0.800\,(0.000)$ & $0.675$ & $0.449$ & $0.872$ & $1228$ \\
Constrained $0.99$   & $0.976\,(0.012)$ & $0.990\,(0.000)$ & $0.777$ & $0.747$ & $0.880$ & $1480$ \\
\bottomrule
\end{tabular}
\end{center}
\end{table}
